\documentclass{amsart}
% For dealing with references we use the comment environment
\usepackage{verbatim}
\newenvironment{reference}{\comment}{\endcomment}
%\newenvironment{reference}{}{}
\newenvironment{slogan}{\comment}{\endcomment}
\newenvironment{history}{\comment}{\endcomment}

% For commutative diagrams we use Xy-pic
\usepackage[all]{xy}

% We use 2cell for 2-commutative diagrams.
\xyoption{2cell}
\UseAllTwocells

% We use multicol for the list of chapters between chapters
\usepackage{multicol}

% This is generally recommended for better output
\usepackage{lmodern}
\usepackage[T1]{fontenc}

% For cross-file-references

% Package for hypertext links:
\usepackage{hyperref}

% For any local file, say "hello.tex" you want to link to please

% Theorem environments.
%
\theoremstyle{plain}
\newtheorem{theorem}[subsection]{Theorem}
\newtheorem{proposition}[subsection]{Proposition}
\newtheorem{lemma}[subsection]{Lemma}

\theoremstyle{definition}
\newtheorem{definition}[subsection]{Definition}
\newtheorem{example}[subsection]{Example}
\newtheorem{exercise}[subsection]{Exercise}
\newtheorem{situation}[subsection]{Situation}

\theoremstyle{remark}
\newtheorem{remark}[subsection]{Remark}
\newtheorem{remarks}[subsection]{Remarks}

\numberwithin{equation}{subsection}

% Macros
%
\def\lim{\mathop{\mathrm{lim}}\nolimits}
\def\colim{\mathop{\mathrm{colim}}\nolimits}
\def\Spec{\mathop{\mathrm{Spec}}}
\def\Hom{\mathop{\mathrm{Hom}}\nolimits}
\def\Ext{\mathop{\mathrm{Ext}}\nolimits}
\def\SheafHom{\mathop{\mathcal{H}\!\mathit{om}}\nolimits}
\def\SheafExt{\mathop{\mathcal{E}\!\mathit{xt}}\nolimits}
\def\Sch{\mathit{Sch}}
\def\Mor{\mathop{\mathrm{Mor}}\nolimits}
\def\Ob{\mathop{\mathrm{Ob}}\nolimits}
\def\Sh{\mathop{\mathit{Sh}}\nolimits}
\def\NL{\mathop{N\!L}\nolimits}
\def\CH{\mathop{\mathrm{CH}}\nolimits}
\def\proetale{{pro\text{-}\acute{e}tale}}
\def\etale{{\acute{e}tale}}
\def\QCoh{\mathit{QCoh}}
\def\Ker{\mathop{\mathrm{Ker}}}
\def\Im{\mathop{\mathrm{Im}}}
\def\Coker{\mathop{\mathrm{Coker}}}
\def\Coim{\mathop{\mathrm{Coim}}}

% Boxtimes
%
\DeclareMathSymbol{\boxtimes}{\mathbin}{AMSa}{"02}

%
% Macros for moduli stacks/spaces
%
\def\QCohstack{\mathcal{QC}\!\mathit{oh}}
\def\Cohstack{\mathcal{C}\!\mathit{oh}}
\def\Spacesstack{\mathcal{S}\!\mathit{paces}}
\def\Quotfunctor{\mathrm{Quot}}
\def\Hilbfunctor{\mathrm{Hilb}}
\def\Curvesstack{\mathcal{C}\!\mathit{urves}}
\def\Polarizedstack{\mathcal{P}\!\mathit{olarized}}
\def\Complexesstack{\mathcal{C}\!\mathit{omplexes}}
% \Pic is the operator that assigns to X its picard group, usage \Pic(X)
% \Picardstack_{X/B} denotes the Picard stack of X over B
% \Picardfunctor_{X/B} denotes the Picard functor of X over B
\def\Pic{\mathop{\mathrm{Pic}}\nolimits}
\def\Picardstack{\mathcal{P}\!\mathit{ic}}
\def\Picardfunctor{\mathrm{Pic}}
\def\Deformationcategory{\mathcal{D}\!\mathit{ef}}

\setlength{\emergencystretch}{3em}
\begin{document}
\title[Stacks Project, Tag 09TG]{392 --- Derived quasi-coherent comparison equivalence from affine etale pushforward compatibility}
\maketitle
\noindent Copyright (C) 2005--2025 Johan de Jong. Stacks Project authors. Source: \href{https://stacks.math.columbia.edu/tag/09TG}{Tag 09TG}.

\noindent Editorial difficulty note (not author text). Difficulty H3: The algebraic-space proof requires more than the scheme support induction: it must lift arbitrary derived morphisms to actual maps after suitable quasi-isomorphisms and control their null homotopies. It proves both fullness and faithfulness using supported cones, then establishes essential surjectivity by an adjunction argument detecting vanishing on every affine etale chart..

\noindent Editorial provenance note (not author text). History: excerpt from revision \texttt{a0444\allowbreak 6e57e\allowbreak c1fbc\allowbreak 252a8\allowbreak 71afc\allowbreak ec775\allowbreak 2fb28\allowbreak 07b14}, spaces-perfect.tex, lines 2272--2469. Reference presentation was adapted; original-author context and core support blocks were retained or added. The mathematical wording is unchanged. Licensed under GNU FDL 1.2 or later; no invariant sections or cover texts. See ../licenses/COPYING. Context blocks reproduce author definitions and supporting results; they are not additional counted problems.

\begin{definition}
\label{definition-elementary-distinguished-square}
Let $S$ be a scheme. A commutative diagram
$$
\xymatrix{
U \times_W V \ar[r] \ar[d] & V \ar[d]^f \\
U \ar[r]^j & W
}
$$
of algebraic spaces over $S$ is called an {\it elementary distinguished square}
if
\begin{enumerate}
\item $U$ is an open subspace of $W$ and $j$ is the inclusion morphism,
\item $f$ is \'etale, and
\item setting $T = W \setminus U$ (with reduced induced
subspace structure) the morphism $f^{-1}(T) \to T$ is an isomorphism.
\end{enumerate}
We will indicate this by saying: ``Let $(U \subset W, f : V \to W)$
be an elementary distinguished square.''
\end{definition}

\noindent
This makes sense by
Properties of Spaces, Lemma
\href{https://stacks.math.columbia.edu/tag/03M1}{lemma properties quasi coherent (Tag 03M1)}
and
Derived Categories, Section \href{https://stacks.math.columbia.edu/tag/06UP}{section triangulated sub (Tag 06UP)}.
Thus we obtain a canonical functor
\begin{equation}
\label{equation-compare}
D(\QCoh(\mathcal{O}_X))
\longrightarrow
D_\QCoh(\mathcal{O}_X)
\end{equation}
see Derived Categories, Equation (\href{https://stacks.math.columbia.edu/tag/06UR}{equation compare (Tag 06UR)}).

\section{The coherator}
\label{section-coherator}

\noindent
Let $S$ be a scheme. Let $X$ be an algebraic space over $S$.
The {\it coherator} is a functor
$$
Q_X :
\textit{Mod}(\mathcal{O}_X)
\longrightarrow
\QCoh(\mathcal{O}_X)
$$
which is right adjoint to the inclusion functor
$\QCoh(\mathcal{O}_X) \to \textit{Mod}(\mathcal{O}_X)$.
It exists for any algebraic space $X$ and moreover the adjunction mapping
$Q_X(\mathcal{F}) \to \mathcal{F}$ is an isomorphism for every
quasi-coherent module $\mathcal{F}$, see
Properties of Spaces, Proposition
\href{https://stacks.math.columbia.edu/tag/077V}{proposition coherator (Tag 077V)}.
Since $Q_X$ is left exact (as a right adjoint) we can consider its
right derived extension
$$
RQ_X :
D(\mathcal{O}_X)
\longrightarrow
D(\QCoh(\mathcal{O}_X)).
$$
Since $Q_X$ is right adjoint to the inclusion functor
$\QCoh(\mathcal{O}_X) \to \textit{Mod}(\mathcal{O}_X)$
we see that $RQ_X$ is right adjoint to the canonical functor
$D(\QCoh(\mathcal{O}_X)) \to D(\mathcal{O}_X)$ by
Derived Categories, Lemma \href{https://stacks.math.columbia.edu/tag/09T5}{lemma derived adjoint functors (Tag 09T5)}.

\medskip\noindent
In this section we will study the functor $RQ_X$. In
Section \href{https://stacks.math.columbia.edu/tag/0CR3}{section better coherator (Tag 0CR3)}
we will study the (closely related) right adjoint to the inclusion functor
$D_\QCoh(\mathcal{O}_X) \to D(\mathcal{O}_X)$ (when it exists).



\begin{lemma}
\label{lemma-argument-proves}
Let $S$ be a scheme. Let $X$ be a quasi-compact and quasi-separated
algebraic space over $S$. Suppose that for every \'etale morphism
$j : V \to W$ with $W \subset X$ quasi-compact open and $V$ affine
the right derived functor
$$
\Phi : D(\QCoh(\mathcal{O}_U)) \to D(\QCoh(\mathcal{O}_W))
$$
of the left exact functor
$j_* : \QCoh(\mathcal{O}_V) \to \QCoh(\mathcal{O}_W)$
fits into a commutative diagram
$$
\xymatrix{
D(\QCoh(\mathcal{O}_V)) \ar[d]_\Phi \ar[r]_{i_V} &
D_\QCoh(\mathcal{O}_V) \ar[d]^{Rj_*} \\
D(\QCoh(\mathcal{O}_W)) \ar[r]^{i_W} &
D_\QCoh(\mathcal{O}_W)
}
$$
Then the functor (\ref{equation-compare})
$$
D(\QCoh(\mathcal{O}_X))
\longrightarrow
D_\QCoh(\mathcal{O}_X)
$$
is an equivalence with quasi-inverse given by $RQ_X$.
\end{lemma}

\begin{proof}
We first use the induction principle to prove $i_X$ is fully faithful.
More precisely, we will use Lemma \href{https://stacks.math.columbia.edu/tag/09IT}{lemma induction principle enlarge (Tag 09IT)}.
Let $(U \subset W, V \to W)$ be an elementary distinguished square
with $V$ affine and $U, W$ quasi-compact open in $X$. Assume that
$i_U$ is fully faithful. We have to show that $i_W$ is fully faithful.
We may replace $X$ by $W$, i.e., we may assume $W = X$ (we do this just
to simplify the notation -- observe that the condition in the
statement of the lemma is preserved under this operation).

\medskip\noindent
Suppose that $A, B$ are objects of $D(\QCoh(\mathcal{O}_X))$.
We want to show that
$$
\Hom_{D(\QCoh(\mathcal{O}_X))}(A, B)
\longrightarrow
\Hom_{D(\mathcal{O}_X)}(i_X(A), i_X(B))
$$
is bijective. Let $T = |X| \setminus |U|$.

\medskip\noindent
Assume first $i_X(B)$ is supported on $T$. In this case the map
$$
i_X(B) \to Rj_{V, *}(i_X(B)|_V) = Rj_{V, *}(i_V(B|_V))
$$
is a quasi-isomorphism
(Lemma \href{https://stacks.math.columbia.edu/tag/08GG}{lemma pushforward with support in open (Tag 08GG)}).
By assumption we have an isomorphism
$i_X(\Phi(B|_V)) \to Rj_{V, *}(i_V(B|_V))$ in $D(\mathcal{O}_X)$.
Moreover, $\Phi$ and ${-}|_V$ are adjoint functors on the derived categories of
quasi-coherent modules (by
Derived Categories, Lemma \href{https://stacks.math.columbia.edu/tag/09T5}{lemma derived adjoint functors (Tag 09T5)}).
The adjunction map $B \to \Phi(B|_V)$ becomes an isomorphism
after applying $i_X$, whence is an isomorphism in
$D(\QCoh(\mathcal{O}_X))$.
Hence
\begin{align*}
\Mor_{D(\QCoh(\mathcal{O}_X))}(A, B)
& =
\Mor_{D(\QCoh(\mathcal{O}_X))}(A, \Phi(B|_V)) \\
& =
\Mor_{D(\QCoh(\mathcal{O}_V))}(A|_V, B|_V) \\
& =
\Mor_{D(\mathcal{O}_V)}(i_V(A|_V), i_V(B|_V)) \\
& =
\Mor_{D(\mathcal{O}_X)}(i_X(A), Rj_{V, *}(i_V(B|_V))) \\
& =
\Mor_{D(\mathcal{O}_X)}(i_X(A), i_X(B))
\end{align*}
as desired. Here we have used that $i_V$ is fully faithful
(Lemma \href{https://stacks.math.columbia.edu/tag/08H0}{lemma affine coherator (Tag 08H0)}).

\medskip\noindent
In general, choose any complex $\mathcal{B}^\bullet$ of quasi-coherent
$\mathcal{O}_X$-modules representing $B$. Next, choose any quasi-isomorphism
$s : \mathcal{B}^\bullet|_U \to \mathcal{C}^\bullet$ of complexes of
quasi-coherent modules on $U$. As $j_U : U \to X$ is
quasi-compact and quasi-separated the functor $j_{U, *}$ transforms
quasi-coherent modules into quasi-coherent modules
(Morphisms of Spaces, Lemma \href{https://stacks.math.columbia.edu/tag/03M9}{lemma pushforward (Tag 03M9)}).
Thus there is a canonical map
$\mathcal{B}^\bullet \to j_{U, *}(\mathcal{B}^\bullet|_U) \to
j_{U, *}\mathcal{C}^\bullet$
of complexes of quasi-coherent modules on $X$.
Set $B'' = j_{U, *}\mathcal{C}^\bullet$ in $D(\QCoh(\mathcal{O}_X))$
and choose a distinguished triangle
$$
B \to B'' \to B' \to B[1]
$$
in $D(\QCoh(\mathcal{O}_X))$. Since the first arrow of the triangle
restricts to an isomorphism over $U$ we see that $B'$ is supported on $T$.
Hence in the diagram
$$
\xymatrix{
\Hom_{D(\QCoh(\mathcal{O}_X))}(A, B'[-1]) \ar[r] \ar[d] &
\Hom_{D(\mathcal{O}_X)}(i_X(A), i_X(B')[-1]) \ar[d] \\
\Hom_{D(\QCoh(\mathcal{O}_X))}(A, B) \ar[r] \ar[d] &
\Hom_{D(\mathcal{O}_X)}(i_X(A), i_X(B)) \ar[d] \\
\Hom_{D(\QCoh(\mathcal{O}_X))}(A, B'') \ar[r] \ar[d] &
\Hom_{D(\mathcal{O}_X)}(i_X(A), i_X(B'')) \ar[d] \\
\Hom_{D(\QCoh(\mathcal{O}_X))}(A, B') \ar[r] &
\Hom_{D(\mathcal{O}_X)}(i_X(A), i_X(B'))
}
$$
we have exact columns and the top and bottom horizontal arrows are
bijective. Finally, choose a complex $\mathcal{A}^\bullet$
of quasi-coherent modules representing $A$.

\medskip\noindent
Let $\alpha : i_X(A) \to i_X(B)$ be a morphism between
in $D(\mathcal{O}_X)$. The restriction $\alpha|_U$ comes from a
morphism in $D(\QCoh(\mathcal{O}_U))$ as $i_U$ is fully faithful.
Hence there exists a choice of
$s : \mathcal{B}^\bullet|_U \to \mathcal{C}^\bullet$ as above
such that $\alpha|_U$ is represented by an actual map of complexes
$\mathcal{A}^\bullet|_U \to \mathcal{C}^\bullet$.
This corresponds to a map of complexes
$\mathcal{A} \to j_{U, *}\mathcal{C}^\bullet$.
In other words, the image of $\alpha$ in
$\Hom_{D(\mathcal{O}_X)}(i_X(A), i_X(B''))$ comes from
an element of $\Hom_{D(\QCoh(\mathcal{O}_X))}(A, B'')$.
A diagram chase then shows that $\alpha$ comes from a morphism
$A \to B$ in $D(\QCoh(\mathcal{O}_X))$. Finally, suppose
that $a : A \to B$ is a morphism of $D(\QCoh(\mathcal{O}_X))$
which becomes zero in $D(\mathcal{O}_X)$. After choosing $\mathcal{B}^\bullet$
suitably, we may assume $a$ is represented by a morphism of complexes
$a^\bullet : \mathcal{A}^\bullet \to \mathcal{B}^\bullet$.
Since $i_U$ is fully faithul the restriction $a^\bullet|_U$ is zero
in $D(\QCoh(\mathcal{O}_U))$. Thus we can choose $s$
such that
$s \circ a^\bullet|_U : \mathcal{A}^\bullet|_U \to \mathcal{C}^\bullet$
is homotopic to zero. Applying the functor $j_{U, *}$ we conclude that
$\mathcal{A}^\bullet \to j_{U, *}\mathcal{C}^\bullet$ is homotopic
to zero. Thus $a$ maps to zero in
$\Hom_{D(\QCoh(\mathcal{O}_X))}(A, B'')$.
Thus we may assume that $a$ is the image of an element
of $b \in \Hom_{D(\QCoh(\mathcal{O}_X))}(A, B'[-1])$.
The image of $b$ in $\Hom_{D(\mathcal{O}_X)}(i_X(A), i_X(B')[-1])$
comes from a $\gamma \in \Hom_{D(\mathcal{O}_X)}(A, B''[-1])$
(as $a$ maps to zero in the group on the right). Since we've
seen above the horizontal arrows are surjective, we see
that $\gamma$ comes from a $c$ in
$\Hom_{D(\QCoh(\mathcal{O}_X))}(A, B''[-1])$
which implies $a = 0$ as desired.

\medskip\noindent
At this point we know that $i_X$ is fully faithful for our original $X$.
Since $RQ_X$ is its right adjoint, we see that
$RQ_X \circ i_X = \text{id}$ (Categories, Lemma
\href{https://stacks.math.columbia.edu/tag/07RB}{lemma adjoint fully faithful (Tag 07RB)}).
To finish the proof we show that for any
$E$ in $D_\QCoh(\mathcal{O}_X)$ the map
$i_X(RQ_X(E)) \to E$ is an isomorphism. Choose a distinguished triangle
$$
i_X(RQ_X(E)) \to E \to E' \to i_X(RQ_X(E))[1]
$$
in $D_\QCoh(\mathcal{O}_X)$. A formal argument using the
above shows that $i_X(RQ_X(E')) = 0$. Thus it suffices to prove that
for $E \in D_\QCoh(\mathcal{O}_X)$ the condition
$i_X(RQ_X(E)) = 0$ implies that $E = 0$. Consider an \'etale morphism
$j : V \to X$ with $V$ affine. By
Lemmas \href{https://stacks.math.columbia.edu/tag/08H0}{lemma affine coherator (Tag 08H0)} and
\href{https://stacks.math.columbia.edu/tag/08GZ}{lemma flat pushforward coherator (Tag 08GZ)}
and our assumption we have
$$
Rj_*(E|_V) = Rj_*(i_V(RQ_V(E|_V))) = i_X(\Phi(RQ_V(E|_V))) =
i_X(RQ_X(Rj_*(E|_V)))
$$
Choose a distinguished triangle
$$
E \to Rj_*(E|_V) \to E' \to E[1]
$$
Apply $RQ_X$ to get a distinguished triangle
$$
0 \to RQ_X(Rj_*(E|_V)) \to RQ_X(E') \to 0[1]
$$
in other words the map in the middle is an isomorphism.
Combined with the string of equalities above we find
that our first distinguished triangle becomes a distinguished triangle
$$
E \to i_X(RQ_X(E')) \to E' \to E[1]
$$
where the middle morphism is the adjunction map. However, the composition
$E \to E'$ is zero, hence $E \to i_X(RQ_X(E'))$ is zero by adjunction!
Since this morphism is isomorphic to the morphism
$E \to Rj_*(E|_V)$ adjoint to $\text{id} : E|_V \to E|_V$ we
conclude that $E|_V$ is zero. Since this holds for all
affine $V$ \'etale over $X$ we conclude $E$ is zero as desired.
\end{proof}
\end{document}
